\chapter{总结与展望}

本章对全文进行总结，并分别列出每个研究内容的贡献创新，并讨论了未来可行的研究工作展望。
\section{全文总结}
随着 GPS 设备与智能移动终端的广泛普及，带有地理位置信息的时空数据呈现出爆炸式增长，这为城市计算与时空数据挖掘提供了丰富的数据基础，也带来了新的研究挑战。特别是在路径规划、接触搜索、联合移动模式挖掘以及轨迹相似性分析等任务中，如何在大规模、动态且存在观测不确定性的环境下进行高效计算，成为学术界和工业界亟需解决的问题。传统方法往往依赖离散轨迹点或静态路径集合，难以满足实时性、可扩展性及准确性要求，尤其在连续到达的行程请求流或大规模移动对象场景下，其性能和适用性存在明显局限。  为应对上述挑战，本文系统研究了面向城市计算的高效算法设计，提出了一系列针对不同场景与任务的算法框架。具体而言，本文涵盖四类核心研究工作：(1) 持续最优路径组合规划：针对动态路网中持续到达的行程请求流，提出了全局协同优化的路径规划方法，通过将时间窗口内的请求作为整体进行联合优化，实现整体通行效率的提升，有效缓解交通拥堵。  (2) 联合移动模式挖掘：面向大规模移动对象轨迹数据，设计了高效的群体模式检测算法，用于高效发现一种宽松的联合移动模式，支持交通管理、拼车推荐及异常行为识别等应用。 (3) 接触搜索与疾病传播分析：针对大规模连续轨迹数据，提出可扩展且低延迟的接触搜索方法，通过运动范围建模有效识别潜在接触事件，为公共卫生监测、疫情防控及风险评估提供数据支撑。  (4) 轨迹相似连接：提出基于运动范围的交集相似性连接框架，突破传统离散点方法的局限，在稀疏采样或不确定条件下仍能稳健刻画对象间潜在交互关系，适用于交通监控、野生动物迁徙分析及接触追踪等场景。  



\subsection{面向大规模行程请求流的持续最优路线组合规划} 
本文在\ref{cha:route search}研究了城市计算应用中面向大规模行程请求流的全局协同路径规划问题。该问题的复杂性在于：行程请求以数据流形式持续到达，早期规划结果会影响未来交通状态，不同行程路径通过共享路段形成复杂耦合关系，从而导致组合搜索空间呈指数级增长。传统单次最优路径规划方法无法应对这种持续性、动态性和全局耦合性的挑战，容易在局部最优决策叠加下形成新的拥堵热点。  

为此，设计了一种高效的自感知批处理解决方案，用于持续优化行程请求流中的路线组合。解决方案包含两个核心步骤：初始路径搜索与批量优化处理。初始路径搜索步骤为每批新到达的行程请求生成高质量的初始路径组合，以避免潜在交通拥堵；批量优化处理步骤则在初始组合基础上进一步优化所有路径，以最小化总通行时间并提升整体路网效率。   研究成果和贡献总结如下：  
\begin{itemize}
    \item 提出了一种新颖的持续最优路线组合问题模型，适用于城市交通流管理、实时路径规划及拥堵预防等多种应用场景。
    \item 设计了高效的自感知批处理算法及剪枝机制，实现对大规模行程请求流的实时优化，兼顾准确性与计算效率。 
    \item 验证了所提方法在真实路网数据上的表现，相较于基线方法展现出更高的性能和可扩展性，在总通行时间优化和实时处理能力上具有显著优势。
\end{itemize}


\subsection{基于轨迹数据的群体联合移动模式挖掘} 
本文在\ref{cha:comovement}研究了城市计算应用中面向大规模移动对象轨迹的联合移动模式挖掘问题。该问题的复杂性在于：轨迹数据规模庞大，移动对象随时间持续变化，联合移动模式需要在空间邻近性、时间持续性及群体规模约束下进行精确判定；同时，传统模式定义在时间上过于严格，导致候选生成与验证成本极高，尤其在实时或近实时应用场景中难以满足计算效率要求。此外，群体成员可能暂时离开群体并重新加入，历史数据与新到达轨迹间存在强动态依赖，使得模式挖掘具有高度动态性与组合耦合性。  

为此，设计了一种针对伴随群体的高效挖掘框架，包含两类核心算法：精确搜索算法（ES-ECMC）和社区感知的近似搜索算法（CS-ACMC）。ES-ECMC算法基于严格的时空约束判定机制，对候选对象对及群体进行精确验证，保证结果完备性与准确性，适用于离线或中小规模数据场景；CS-ACMC算法通过近似计算与结构化聚类方法压缩搜索空间，实现高效增量挖掘，满足实时或近实时应用需求。整体解决方案的步骤为：首先在近似联合移动对搜索阶段，通过网格划分将轨迹离散化，并快速估计对象间联合移动程度生成候选对；随后在社区感知群体发现阶段，将候选对构建为图结构并利用社区检测算法进行聚类，从而高效识别伴随群体。 研究成果和贡献总结如下：
\begin{itemize}
    \item 提出了伴随群体模式定义，兼顾空间邻近性、时间持续性与群体规模约束，并适应动态轨迹数据环境，具备广泛应用价值，包括交通监测、疫情防控及公共安全预警等场景。
    \item 设计了精确与近似两类算法（ES-ECMC与CS-ACMC），有效解决了准确性与效率之间的权衡问题。
    \item 验证了所提方法在用户可交互时间内高效发现伴随群体的能力，同时展示了在大规模数据环境下的可扩展性与实用性；
\end{itemize}  
总体而言，本章在理论与实践层面系统研究了伴随群体发现问题，提出了高
效、可扩展且具有应用价值的解决方案，为大规模动态轨迹数据环境下的群体行
为挖掘提供了新的研究方法与技术工具。


\subsection{大规模移动对象上的持续密切接触发现} 
本文在\ref{cha:contact}研究了城市计算应用中面向大规模动态移动对象的实时密切接触检测问题。该问题的复杂性在于：轨迹数据以持续流的形式到达，移动对象数量巨大且不断变化，同时需要在空间邻近性与时间持续性约束下实时识别潜在密切接触对象。此外，密切接触具有递归扩展特性，即新发现的接触对象将继续参与后续检测，导致传播链持续增长，增加了算法的计算复杂度与状态维护难度。在突发公共卫生事件或大型人群聚集场景下，若无法高效处理轨迹流，将难以及时识别关键接触节点，从而影响风险控制决策。  

为此，设计了一种针对持续密切接触搜索问题的高效解决方案。解决方案包含两类核心算法：精确实时搜索算法和近似实时搜索算法。精确实时搜索算法通过严格的空间与时间约束判定，对候选对象进行精确验证，保证检测结果的完备性与准确性；近似实时搜索算法在保证不漏检的前提下，通过首尾优先检查、跳跃检测及近似机制显著降低计算量，实现高效实时处理。整体解决方案的步骤为：首先对轨迹对象进行分组过滤，减少候选对象数量；随后基于距离上下界的预检查策略避免不必要的精确距离计算；最后在滑动时间窗口内维护对象间接触状态，实现持续更新与递归扩展。  研究成果和贡献总结如下： 
\begin{itemize}
    \item 定义了持续密切接触搜索问题，为流行病传播监测、公共卫生预警及风险追踪提供统一的数据处理框架，明确了递归扩展与实时更新特性。
    \item 设计了精确与近似两类实时处理算法，结合分组过滤、预检查及跳跃检测等优化策略，有效解决了准确性与效率之间的权衡难题。
    \item 验证了所提出方法具有良好的可扩展性与稳定性：在保证检测精度的同时，实现了大规模动态场景下的高效实时密切接触检测。
\end{itemize} 
总体而言，本章系统研究了持续密切接触搜索问题，提出了高效、可扩展且具
有应用价值的解决方案，为大规模流式轨迹数据环境下的实时疫情监控与风险管
理提供了新的技术方法与实践参考。

\subsection{面向时空运动范围的轨迹相似度连接} 
本文在\ref{cha:join}揭示了城市计算应用中移动对象相似性连接问题中的一个关键挑战——在稀疏采样条件下如何准确刻画对象间潜在交互关系。传统基于离散采样点的相似性度量方法在对象轨迹采样频率不足或存在观测缺失时容易低估真实相似性，导致潜在对象对被误剪枝。为解决这一问题，提出了交集相似性连接问题，通过引入“运动范围”概念，将对象在相邻观测样本之间可能到达的空间区域建模为运动范围，并基于这些区域的交集程度进行相似性度量，从而在稀疏采样或数据噪声条件下更稳健地刻画潜在交互。

交集相似性连接问题在多个实际场景中具有重要应用价值。例如，在城市交通监测中，车辆轨迹的运动范围可映射为道路子路段或交叉口集合，通过交集相似性识别潜在交通关联，有助于拥堵预测与异常路段识别；在野生动物迁徙分析中，运动范围结合栖息地与迁徙通道信息，可识别共享走廊或高频活动区域，实现低采样频率下的关键交互检测；在公共卫生领域，个体潜在接触事件可通过运动范围交集推断，从而支持疾病传播建模与风险评估。不同应用虽对象类型与特征不同，但核心需求均为在观测不确定条件下识别潜在运动交集关系，该研究提供了统一且可扩展的解决框架。

在建模与计算上具有以下关键创新：首先，明确将相似性度量从离散点级距离扩展到连续时间下的区域交集，允许在理论上处理未观测时间段的潜在交互；其次，将连续时间区间内的交集相似性定义为时间平均函数，既保留连续性信息，又便于离散化近似计算。该建模显著提升了稀疏采样条件下潜在对象对识别能力，同时引入连续时间处理与大规模候选对计算的两大挑战：一是连续时间计算的复杂性，需设计可计算的离散化或可控近似方法；二是高效连接的挑战，现有点或静态区域相似性算法无法直接支持连续时间运动范围的交集计算，否则候选对象对数量会急剧膨胀。

为应对上述挑战，提出了基于时间区间运动范围的高效求解框架。核心思想是通过混合 球树索引对对象运动范围进行层次化组织，并结合重划分策略与多级剪枝技术，有效过滤候选对象对，从而在保证正确性与完整性的前提下显著降低计算开销。该框架适用于稀疏采样轨迹数据，能够高效识别满足相似性阈值的对象对，实现连续时间区间下的交集相似性计算。研究成果和贡献总结如下：  
\begin{itemize}
    \item 提出了交集相似性连接问题，将相似性度量从离散点扩展到连续时间下的运动范围交集，显著增强稀疏采样轨迹潜在交互建模能力。
    \item 构建了混合球树层次索引结构，并结合重划分与多级剪枝策略，实现大规模连续时间区间下的高效连接计算。
    \item 验证了所提出方法的可扩展性与实用价值：结果表明方法在保证准确性的前提下，相较于现有基线方法平均运行时间降低约 3 倍。
\end{itemize}

总之，本文针对稀疏采样轨迹数据中潜在运动交集关系的识别提出了统一建模与高效计算框架，为交通分析、生态监测及公共卫生等多种实际场景提供了可行且高效的技术方案。

\section{展望}
本文针对城市计算应用中的4个需求，即面向大规模行程请求流的持续最优路线组合规划需求，基于轨迹数据的群体联合移动模式挖掘需求，面向大规模移动对象的持续密切接触查询需求，以及面向时空运动范围的轨迹相似度连接需求展开了深入研究，发表了一系列研究成果，但仍存在一些值得进一步研究和探索的空间。随着移动对象数据规模和实时更新频率的持续增长，单机处理在计算能力与内存资源上面临显著瓶颈。未来可能的进一步研究方向可以向分布式环境、深度学习优化、时空大模型拓展，潜在研究方向包括：

\paragraph{分布式查询}
针对大规模移动对象或位置标注数据，当数据量超出单机处理能力时，可在分布式框架中研究并行化的索引机制和查询算法。研究重点可包括：面向海量时空-文本数据的并行查询架构，优化任务调度与负载均衡；以及分布式数据迁移与副本管理策略，以保证连续查询过程中的数据一致性与高可用性。在分布式环境下，将交集相似性连接或连续密切接触搜索等计算任务进行分片处理，同时保持跨节点候选对象的完整性与准确性，是提高系统可扩展性与实时性的重要方向。相关研究可包括分布式滑动窗口维护、跨节点运动范围交集计算以及并行剪枝策略设计。分布式处理环境中，节点故障、网络延迟及数据倾斜可能导致查询性能下降。未来可探索自适应任务调度、容错恢复机制以及数据划分优化方法，以保障大规模连续时空查询的稳定性与效率。

\paragraph{基于深度学习的方法优化}
随着移动对象数据的持续增长与复杂化，传统基于规则或索引的算法在处理大规模、动态、多源数据时面临性能和可扩展性瓶颈。深度学习技术因其强大的表示学习和模式捕捉能力，为未来时空数据分析与连续查询提供了新的研究契机。具体而言，潜在的研究方向可包括以下几个方面：（1）在大规模行程请求流场景下，如何为海量移动对象实时生成最优路线组合是关键问题。深度强化学习可用于建模连续决策过程，将路线选择、时间窗口约束和交通动态因素作为状态与奖励函数，实现高效、可扩展的路径优化。通过结合历史轨迹数据与实时交通信息，模型可学习群体行程模式，从而在动态环境中提供持续优化的路线方案。（2）群体行为分析和联合移动模式识别在城市规划、公共交通调度及疫情传播监控中具有重要意义。深度学习方法，如图神经网络或时空卷积网络，可以将移动对象及其轨迹视为时空图结构或序列数据，从而自动捕捉潜在的群体协同行为、同步移动模式及关键影响节点，实现对复杂群体行为的高效挖掘。（3）持续密切接触检测涉及大规模对象在动态环境下的空间邻近与时间连续性判定。深度学习模型可通过学习对象轨迹的潜在运动分布和邻近概率，实现对未观测时间段潜在接触的预测，从而在海量数据流环境下支持实时、鲁棒的密切接触检测。结合序列建模与空间嵌入方法，可以有效提升接触预测的精度与响应速度。（4） 在稀疏采样轨迹或不完整数据条件下，基于传统规则的运动范围交集计算容易受限。深度表示学习可通过将轨迹映射到低维时空潜在空间，捕捉连续时间区间内的潜在重叠关系，从而高效完成交集相似性连接。此外，生成式模型可以辅助补全未观测轨迹，进一步增强相似性连接的准确性与鲁棒性。

\paragraph{时空大模型设计}
未来可以构建一个统一的时空大模型框架，将轨迹数据分析、群体联合移动模式挖掘、持续密切接触查询以及轨迹相似度连接等任务整合在同一平台中，并提出接口用于用户交互。该框架可通过深度学习方法实现统一时空表示，将多源轨迹、行为和环境信息映射到低维潜在空间，从而支持个体轨迹预测、群体模式识别、实时接触检测和交集相似性计算。针对大规模和高频数据，可在分布式计算环境下进行流式处理与增量模型更新，实现高效并行计算与实时响应。同时，框架可提供统一接口与查询语义，便于不同应用场景快速接入，包括智能交通、公共卫生监控和生态保护等。结合可解释性和隐私保护技术，系统既能输出可理解的决策依据，又保障个体数据安全。整体来看，该统一时空大模型有望通过深度学习和大数据优化，实现跨任务、跨数据源的高效、实时和可扩展时空智能分析。